\section*{Tóm tắt}
\thispagestyle{empty}

Nghiên cứu này tập trung phát triển hệ thống gợi ý tuần tự đa phương thức ứng dụng trên tập dữ liệu Amazon Reviews, nhằm tối ưu hóa trải nghiệm tìm kiếm và mua sắm quà tặng trực tuyến. Mục tiêu cốt lõi của hệ thống là hỗ trợ người dùng tìm kiếm sản phẩm phù hợp nhất thông qua việc phân tích sâu nội dung đa phương thức bao gồm văn bản và hình ảnh, kết hợp với chuỗi hành vi và sở thích lịch sử để đưa ra các gợi ý có tính cá nhân hóa cao.\\

Về mặt phương pháp, hệ thống được xây dựng dựa trên quy trình truy xuất và xếp hạng đa tầng (Retrieve-and-Rerank). Khung kiến trúc nền tảng bao gồm mô hình Tháp kép (Two-Tower) kết hợp cùng mạng Cross-Encoder. Trong pha truy xuất sơ bộ, tháp sản phẩm ứng dụng cơ chế cổng hợp nhất để tổng hợp đặc trưng văn bản từ BERT, hình ảnh từ ViT và dữ liệu bảng. Đồng thời, tháp người dùng mô hình hóa chuỗi tương tác tuần tự bằng mạng Transformer cải tiến từ SASRec, tích hợp khả năng nhận biết cảm xúc để học hỏi từ các phản hồi tích cực và tiêu cực. Để tối ưu hóa hiệu suất và tăng tính minh bạch, Mô hình Ngôn ngữ Lớn được ứng dụng sâu vào hệ thống với hai vai trò cốt lõi: tự động sinh dữ liệu truy vấn từ đánh giá lịch sử để cung cấp đầu vào huấn luyện cho Cross-Encoder, và tạo ra các lời giải thích tự nhiên cho kết quả gợi ý.\\

Quá trình suy luận diễn ra qua một quy trình ba giai đoạn chặt chẽ: thứ nhất, truy xuất tập ứng viên sơ bộ dựa trên độ tương đồng ngữ nghĩa trong không gian vector tiềm ẩn; thứ hai, xếp hạng lại thông qua mạng Cross-Encoder nhằm tính toán trực tiếp mức độ giao thoa giữa ý định tìm kiếm của người dùng và từng sản phẩm cụ thể; và cuối cùng, tinh chỉnh thứ hạng chuyên sâu kết hợp sinh lời giải thích minh bạch bằng LLM. Tại bước cuối cùng, LLM đóng vai trò đánh giá ngữ cảnh và suy luận mối liên hệ giữa nhu cầu hiện tại với lịch sử hành vi để chọn lọc danh sách kết quả tối ưu. Ngoài ra, vấn đề khởi động lạnh đối với người dùng mới được giải quyết thông qua thuật toán đánh giá có trọng số. Kết quả thực nghiệm cho thấy mô hình không chỉ nắm bắt hiệu quả các mối quan hệ ngữ nghĩa đa phương thức và diễn biến sở thích người dùng, mà còn tận dụng tốt khả năng suy luận của LLM để nâng cao chất lượng gợi ý cuối cùng.\\
\clearpage
\pagenumbering{arabic}
